% =====================================================================
%  TD : CONTINUITÉ DES FONCTIONS — 15 exercices progressifs
%  Terminale — Spécialité Mathématiques
%  Lycée Français International Jean CHARCOT — EL JADIDA
%  Année scolaire 2026 - 2027
%  Compilation : pdflatex (deux passes)
% =====================================================================
\documentclass[11pt,a4paper]{article}

% ---------------------------------------------------------------------
%  Encodage, langue, polices
% ---------------------------------------------------------------------
\usepackage[utf8]{inputenc}
\usepackage[T1]{fontenc}
\usepackage{lmodern}
\usepackage[french]{babel}
\usepackage{microtype}

% ---------------------------------------------------------------------
%  Mathématiques
% ---------------------------------------------------------------------
\usepackage{amsmath,amssymb,mathtools}
\usepackage{mathrsfs}

% ---------------------------------------------------------------------
%  Mise en page
% ---------------------------------------------------------------------
\usepackage[a4paper,top=2.3cm,bottom=2.3cm,left=2cm,right=2cm,
            headheight=15pt]{geometry}
\usepackage[table]{xcolor}
\usepackage{array}
\usepackage{multicol}
\usepackage{enumitem}
\usepackage{fancyhdr}
\usepackage{titlesec}
\usepackage{lastpage}
\usepackage{needspace}
\usepackage{setspace}
\usepackage{pifont}
\usepackage{listings}
\usepackage[most]{tcolorbox}

% ---------------------------------------------------------------------
%  Figures
% ---------------------------------------------------------------------
\usepackage{tikz}
\usetikzlibrary{arrows.meta,calc,babel}
\usepackage{tkz-tab}

\usepackage[hidelinks]{hyperref}

% ---------------------------------------------------------------------
%  Interlignage aéré
% ---------------------------------------------------------------------
\setstretch{1.25}
\setlength{\parindent}{0pt}
\setlength{\parskip}{6pt plus 2pt minus 1pt}
\setlength{\abovedisplayskip}{8pt plus 2pt}
\setlength{\belowdisplayskip}{8pt plus 2pt}

% Ne jamais isoler une ligne en haut ou en bas de page
\widowpenalty=10000
\clubpenalty=10000
\displaywidowpenalty=10000

% Remplissage des pages : aucune coupure « favorisée ».
% \needspace{h} passe à la page suivante seulement si l'espace restant est < h.
\makeatletter
\renewcommand{\needspace}[1]{\par\begingroup\setlength{\dimen@}{#1}%
  \dimen@ii\pagegoal\advance\dimen@ii-\pagetotal
  \ifdim\dimen@>\dimen@ii\ifdim\dimen@ii>\z@\vfil\fi\break\fi\endgroup}
\@secpenalty=0
\@beginparpenalty=0
\@endparpenalty=0
\@itempenalty=0
\makeatother

% ---------------------------------------------------------------------
%  Couleurs
% ---------------------------------------------------------------------
\definecolor{bleu}{RGB}{20,55,120}
\definecolor{bleuclair}{RGB}{232,240,252}
\definecolor{rouge}{RGB}{170,30,45}
\definecolor{vert}{RGB}{25,115,65}
\definecolor{vertclair}{RGB}{233,247,238}
\definecolor{orange}{RGB}{200,100,10}
\definecolor{orangeclair}{RGB}{254,243,228}
\definecolor{gris}{RGB}{90,90,90}
\definecolor{grisclair}{RGB}{244,244,244}
\definecolor{or}{RGB}{215,160,20}

% ---------------------------------------------------------------------
%  Titres des parties
% ---------------------------------------------------------------------
\renewcommand{\thesection}{\Alph{section}}
\newcommand{\boitesection}[1]{\colorbox{bleu}{\parbox{\dimexpr\linewidth-2\fboxsep}{%
  \strut Partie~\thesection~--~#1}}}
\titleformat{\section}
  {\needspace{10\baselineskip}\normalfont\Large\bfseries\color{white}}
  {}{0pt}{\boitesection}
\titlespacing*{\section}{0pt}{16pt}{10pt}

% ---------------------------------------------------------------------
%  En-têtes et pieds de page
% ---------------------------------------------------------------------
\pagestyle{fancy}
\fancyhf{}
\fancyhead[L]{\small\color{gris}TD -- Continuité des fonctions}
\fancyhead[R]{\small\color{gris}TSpeMath}
\fancyfoot[L]{\small\color{gris}LFI Jean CHARCOT -- EL JADIDA}
\fancyfoot[C]{\small\color{gris}\thepage/\pageref{LastPage}}
\fancyfoot[R]{\small\color{gris}K. EL GHAZOULI}
\renewcommand{\headrulewidth}{0.4pt}
\renewcommand{\footrulewidth}{0.4pt}

% ---------------------------------------------------------------------
%  Niveaux de difficulté
% ---------------------------------------------------------------------
\newcommand{\etoile}{\textcolor{or}{\ding{72}}}
\newcommand{\etoilevide}{\textcolor{or}{\ding{73}}}
\newcommand{\niveauA}{\etoile\etoilevide\etoilevide}
\newcommand{\niveauB}{\etoile\etoile\etoilevide}
\newcommand{\niveauC}{\etoile\etoile\etoile}

% ---------------------------------------------------------------------
%  Environnement « exercice »
% ---------------------------------------------------------------------
\newtcolorbox[auto counter]{exercice}[2][]{
  enhanced, breakable, colback=white, colframe=bleu, boxrule=0.8pt,
  arc=2mm, sharp corners=south, left=4mm, right=4mm, top=2.5mm, bottom=2.5mm,
  colbacktitle=bleuclair, coltitle=bleu, fonttitle=\bfseries,
  toptitle=1.2mm, bottomtitle=1.2mm,
  title={Exercice~\thetcbcounter\ifx&#1&\else~--~#1\fi\hfill
         {\normalfont\small Niveau}~#2},
  before upper pre={\setlength{\parskip}{6pt}},
  before pre={\par\needspace{8\baselineskip}},
  before skip=10pt, after skip=10pt}

\newtcolorbox{indication}{
  enhanced, colback=grisclair, colframe=gris, boxrule=0pt, leftrule=1mm,
  sharp corners, left=3mm, right=2mm, top=1.5mm, bottom=1.5mm,
  title={Indication}, attach title to upper, after title={\ :\ },
  coltitle=gris, fonttitle=\bfseries\itshape, fontupper=\itshape,
  before skip=6pt, after skip=4pt}

% Questions : 1. 2. 3. puis a) b) c)
\setlist[enumerate,1]{label=\textcolor{bleu}{\textbf{\arabic*.}}, leftmargin=*,
  itemsep=4pt, topsep=3pt, parsep=2pt}
\setlist[enumerate,2]{label=\textcolor{bleu}{\alph*)}, leftmargin=*,
  itemsep=4pt, topsep=3pt}
\setlist[itemize]{label=\textcolor{bleu}{$\bullet$}, leftmargin=*, itemsep=4pt,
  topsep=4pt}

% ---------------------------------------------------------------------
%  Macros mathématiques
% ---------------------------------------------------------------------
\newcommand{\R}{\mathbb{R}}
\newcommand{\N}{\mathbb{N}}
\newcommand{\Z}{\mathbb{Z}}
\newcommand{\Cf}{\mathscr{C}_f}
\newcommand{\e}{\mathrm{e}}
\newcommand{\intff}[2]{\mathopen{[}#1\,;\,#2\mathclose{]}}
\newcommand{\intoo}[2]{\mathopen{]}#1\,;\,#2\mathclose{[}}
\newcommand{\intfo}[2]{\mathopen{[}#1\,;\,#2\mathclose{[}}
\newcommand{\intof}[2]{\mathopen{]}#1\,;\,#2\mathclose{]}}
\newcommand{\limite}[2]{\lim\limits_{#1}#2}
\newcommand{\limg}[2]{\lim\limits_{x\to #1^-}#2}
\newcommand{\limd}[2]{\lim\limits_{x\to #1^+}#2}

% ---------------------------------------------------------------------
%  Code Python
% ---------------------------------------------------------------------
\lstset{
  language=Python, basicstyle=\ttfamily\small\linespread{1}\selectfont,
  keywordstyle=\color{bleu}\bfseries, commentstyle=\color{vert}\itshape,
  numbers=left, numberstyle=\tiny\color{gris}, frame=single,
  rulecolor=\color{gris!50}, backgroundcolor=\color{grisclair},
  showstringspaces=false, columns=fullflexible, keepspaces=true,
  literate={é}{{\'e}}1 {è}{{\`e}}1 {à}{{\`a}}1 {…}{{\dots}}1
}

% =====================================================================
\begin{document}
% =====================================================================

% ---------------------------------------------------------------------
%  EN-TÊTE DE LA FICHE
% ---------------------------------------------------------------------
\thispagestyle{fancy}
\begin{tcolorbox}[enhanced, colback=white, colframe=bleu, boxrule=1.2pt,
  arc=3mm, left=5mm, right=5mm, top=3mm, bottom=3mm, halign=center]
\setstretch{1.15}
{\large\bfseries\color{bleu} Lycée Français International Jean CHARCOT - EL JADIDA\par}
\vspace{1mm}
{Année scolaire 2026 - 2027\par}
{\bfseries Niveau : Terminale Spécialité Mathématiques\par}
\vspace{2mm}
{\color{bleu}\rule{0.7\linewidth}{0.8pt}\par}
\vspace{2mm}
{\fontsize{24}{30}\selectfont\bfseries\color{bleu}\spaceskip=0.4em
 TD -- Continuité des fonctions\par}
\vspace{1mm}
{\itshape 15 exercices progressifs\par}
\end{tcolorbox}

\begin{tcolorbox}[enhanced, colback=orangeclair, colframe=orange, boxrule=0.6pt,
  arc=2mm, left=4mm, right=4mm, top=2mm, bottom=2mm]
\textbf{Niveaux :}\hfill \niveauA\ application directe du cours
\hfill \niveauB\ approfondissement \hfill \niveauC\ type bac

\textbf{Rappel :} $f$ est continue en $a$ si et seulement si
$\limg{a}{f(x)}=\limd{a}{f(x)}=f(a)$.
\end{tcolorbox}

% =====================================================================
\section{Continuité en un point}
% =====================================================================

% ----------------------------------------------------------- Ex 1
\begin{exercice}[Lecture graphique]{\niveauA}
\begin{minipage}[c]{0.5\linewidth}
On donne ci-contre la courbe représentative d'une fonction $f$ définie sur
$\intff{-3}{4}$.

Un point plein \tikz\fill (0,0) circle (2pt); appartient à la courbe ;
un point creux \tikz\draw[fill=white] (0,0) circle (2pt); n'appartient pas
à la courbe.
\end{minipage}\hfill
\begin{minipage}[c]{0.46\linewidth}
\centering
\begin{tikzpicture}[scale=0.72]
  \draw[gris!30] (-3.4,-0.6) grid (4.4,3.6);
  \draw[-{Stealth}] (-3.5,0) -- (4.6,0) node[right] {$x$};
  \draw[-{Stealth}] (0,-0.7) -- (0,3.8) node[above] {$y$};
  \foreach \k in {-3,-2,-1,1,2,3,4}{\node[below, font=\scriptsize] at (\k,0) {$\k$};}
  \foreach \k in {1,2,3}{\node[left, font=\scriptsize] at (0,\k) {$\k$};}
  \node[below left, font=\scriptsize] at (0,0) {$O$};
  \draw[very thick, rouge] (-3,0) -- (-1,2);
  \draw[very thick, rouge, domain=-1:1, samples=40] plot (\x,{\x*\x});
  \draw[very thick, rouge] (1,2) -- (2,1) -- (4,3);
  \fill[rouge] (-3,0) circle (2.4pt);
  \draw[thick, rouge, fill=white] (-1,2) circle (2.4pt);
  \fill[rouge] (-1,1) circle (2.4pt);
  \fill[rouge] (1,1) circle (2.4pt);
  \draw[thick, rouge, fill=white] (1,2) circle (2.4pt);
  \fill[rouge] (2,1) circle (2.4pt);
  \fill[rouge] (4,3) circle (2.4pt);
\end{tikzpicture}
\end{minipage}

\begin{enumerate}
  \item Lire $f(-1)$, $f(1)$ et $f(2)$.
  \item Déterminer graphiquement les limites suivantes :
  \[
    \limg{-1}{f(x)},\quad \limd{-1}{f(x)},\quad
    \limg{1}{f(x)},\quad \limd{1}{f(x)},\quad
    \limg{2}{f(x)},\quad \limd{2}{f(x)}.
  \]
  \item La fonction $f$ est-elle continue en $-1$ ? en $1$ ? en $2$ ? Justifier.
  \item Donner les plus grands intervalles sur lesquels $f$ est continue.
\end{enumerate}
\end{exercice}

% ----------------------------------------------------------- Ex 2
\begin{exercice}[Continuité en un point par le calcul]{\niveauA}
Dans chaque cas, étudier la continuité de la fonction au point $a$ indiqué,
en calculant l'image, la limite à gauche et la limite à droite.
\begin{multicols}{2}
\begin{enumerate}
  \item $f(x)=\begin{cases} x^2-1 & \text{si } x<2\\ 2x-1 & \text{si } x\geqslant 2\end{cases}$
        \quad en $a=2$.
  \item $g(x)=\begin{cases} x+1 & \text{si } x\leqslant 0\\ x^2+2 & \text{si } x>0\end{cases}$
        \quad en $a=0$.
  \item $h(x)=\begin{cases} \e^{x} & \text{si } x<0\\ 1-x & \text{si } x\geqslant 0\end{cases}$
        \quad en $a=0$.
  \item $k(x)=\begin{cases} \dfrac{x^2-4}{x-2} & \text{si } x\neq 2\\[6pt] 4 & \text{si } x=2\end{cases}$
        \quad en $a=2$.
\end{enumerate}
\end{multicols}
\begin{indication}
Pour la question 4, factoriser le numérateur.
\end{indication}
\end{exercice}

% ----------------------------------------------------------- Ex 3
\begin{exercice}[Recherche d'un paramètre]{\niveauA}
\begin{enumerate}
  \item Soit $m$ un réel et $f$ la fonction définie sur $\R$ par
  \[
    f(x)=\begin{cases} \e^{x}+m & \text{si } x<0,\\ 2x+3 & \text{si } x\geqslant 0.\end{cases}
  \]
  Déterminer la valeur de $m$ pour laquelle $f$ est continue sur $\R$.
  \item Soit $a$ un réel et $g$ la fonction définie sur $\R$ par
  \[
    g(x)=\begin{cases} \dfrac{x^2-1}{x-1} & \text{si } x\neq 1,\\[6pt] a & \text{si } x=1.\end{cases}
  \]
  Déterminer la valeur de $a$ pour laquelle $g$ est continue en $1$.
\end{enumerate}
\end{exercice}

% ----------------------------------------------------------- Ex 4
\begin{exercice}[Deux paramètres]{\niveauB}
Soit $a$ et $b$ deux réels et $f$ la fonction définie sur $\R$ par
\[
  f(x)=\begin{cases}
    2x+a & \text{si } x<-1,\\
    x^2 & \text{si } -1\leqslant x\leqslant 2,\\
    bx-2 & \text{si } x>2.
  \end{cases}
\]
\begin{enumerate}
  \item Justifier que $f$ est continue sur $\intoo{-\infty}{-1}$, sur
        $\intoo{-1}{2}$ et sur $\intoo{2}{+\infty}$.
  \item Déterminer $a$ pour que $f$ soit continue en $-1$.
  \item Déterminer $b$ pour que $f$ soit continue en $2$.
  \item Tracer l'allure de la courbe de $f$ pour les valeurs trouvées.
\end{enumerate}
\end{exercice}

% ----------------------------------------------------------- Ex 5
\begin{exercice}[Autour de la partie entière]{\niveauB}
On rappelle que, pour tout réel $x$, $E(x)$ désigne l'unique entier relatif $n$
tel que $n\leqslant x<n+1$.

On considère la fonction $f$ définie sur $\R$ par $f(x)=x-E(x)$.
\begin{enumerate}
  \item Calculer $f(0{,}3)$, $f(2)$, $f(2{,}75)$ et $f(-1{,}5)$.
  \item Exprimer $f(x)$ en fonction de $x$ sur $\intfo{0}{1}$, puis sur $\intfo{1}{2}$.
  \item Tracer la courbe de $f$ sur $\intfo{-2}{3}$.
  \item Calculer $\limg{1}{f(x)}$ et $\limd{1}{f(x)}$. La fonction $f$ est-elle
        continue en $1$ ?
  \item Sur quels intervalles la fonction $f$ est-elle continue ?
  \item \textbf{Pour aller plus loin.} On pose $h(x)=E(x)+\sqrt{x-E(x)}$.
        Montrer que $h$ est continue en $1$.
\end{enumerate}
\end{exercice}

% =====================================================================
\section{Opérations et dérivabilité}
% =====================================================================

% ----------------------------------------------------------- Ex 6
\begin{exercice}[Justifier la continuité]{\niveauA}
Justifier que chacune des fonctions suivantes est continue sur l'intervalle
(ou les intervalles) indiqué(s).
\begin{enumerate}
  \item $f(x)=3x^4-2x+5$ \quad sur $\R$.
  \item $g(x)=\dfrac{x+1}{x^2-4}$ \quad sur $\intoo{-\infty}{-2}$, $\intoo{-2}{2}$
        et $\intoo{2}{+\infty}$.
  \item $h(x)=\sqrt{2x-6}$ \quad sur $\intfo{3}{+\infty}$.
  \item $k(x)=x\,\e^{-x}+\cos(3x)$ \quad sur $\R$.
  \item $m(x)=\sqrt{x^2+x+1}$ \quad sur $\R$.
  \item $n(x)=\dfrac{1}{1+\e^{x}}$ \quad sur $\R$.
\end{enumerate}
\begin{indication}
Pour la question 5, montrer que $x^2+x+1>0$ pour tout réel $x$ (discriminant).
\end{indication}
\end{exercice}

% ----------------------------------------------------------- Ex 7
\begin{exercice}[Continuité et dérivabilité]{\niveauB}
On considère les fonctions $f$ et $g$ définies sur $\R$ par
\[
  f(x)=\begin{cases} x^2 & \text{si } x\leqslant 1\\ 2x-1 & \text{si } x>1\end{cases}
  \qquad\text{et}\qquad
  g(x)=\begin{cases} x^2 & \text{si } x\leqslant 1\\ x & \text{si } x>1.\end{cases}
\]
\begin{enumerate}
  \item Montrer que $f$ et $g$ sont continues en $1$.
  \item Pour $h\neq 0$, calculer le taux de variation $\dfrac{f(1+h)-f(1)}{h}$ pour
        $h<0$, puis pour $h>0$. La fonction $f$ est-elle dérivable en $1$ ?
  \item Même question pour la fonction $g$.
  \item Interpréter graphiquement les résultats des questions 2 et 3.
  \item \textbf{Vrai ou faux ?} Justifier chaque réponse.
  \begin{enumerate}
    \item Si une fonction est continue en $a$, alors elle est dérivable en $a$.
    \item Si une fonction n'est pas continue en $a$, alors elle n'est pas
          dérivable en $a$.
  \end{enumerate}
\end{enumerate}
\end{exercice}

% =====================================================================
\section{Théorème des valeurs intermédiaires}
% =====================================================================

% ----------------------------------------------------------- Ex 8
\begin{exercice}[Existence de solutions]{\niveauA}
Soit $f$ la fonction définie sur $\R$ par $f(x)=x^3-4x+1$.
\begin{enumerate}
  \item Montrer que l'équation $f(x)=0$ admet au moins une solution dans
        $\intff{0}{1}$.
  \item Calculer $f(-3)$, $f(-2)$ et $f(2)$. En déduire que l'équation
        $f(x)=0$ admet au moins trois solutions dans $\R$.
  \item On considère la fonction inverse $u:x\mapsto\dfrac{1}{x}$. On a
        $u(-1)<0$ et $u(1)>0$.

        Peut-on en déduire que l'équation $\dfrac{1}{x}=0$ admet une solution
        dans $\intff{-1}{1}$ ? Expliquer.
\end{enumerate}
\end{exercice}

% ----------------------------------------------------------- Ex 9
\begin{exercice}[Lecture d'un tableau de variations]{\niveauA}
\begin{minipage}{\linewidth}
\setlength{\parskip}{6pt}
Soit $f$ une fonction continue sur $\intff{-4}{5}$ dont le tableau de variations
est le suivant :
\begin{center}
\setstretch{1}
\begin{tikzpicture}
\tkzTabInit[lgt=1.6, espcl=2.6]{$x$/0.8, $f(x)$/1.6}{$-4$, $-1$, $3$, $5$}
\tkzTabVar{+/$2$, -/$-3$, +/$4$, -/$1$}
\end{tikzpicture}
\end{center}
\end{minipage}
\begin{enumerate}
  \item Déterminer le nombre de solutions de chacune des équations suivantes :
  \[
    \text{a) } f(x)=0 \qquad \text{b) } f(x)=3 \qquad \text{c) } f(x)=-4
    \qquad \text{d) } f(x)=1{,}5 .
  \]
  \item Rédiger soigneusement la justification pour l'équation $f(x)=0$.
  \item Donner le signe de $f(x)$ selon les valeurs de $x$ (on notera $\alpha$
        et $\beta$ les solutions de $f(x)=0$).
\end{enumerate}
\end{exercice}

% ----------------------------------------------------------- Ex 10
\begin{exercice}[Existence, unicité, encadrement]{\niveauB}
Soit $f$ la fonction définie sur $\R$ par $f(x)=\e^{x}+x$.
\begin{enumerate}
  \item Déterminer les limites de $f$ en $-\infty$ et en $+\infty$.
  \item Étudier les variations de $f$ sur $\R$ et dresser son tableau de variations.
  \item Montrer que l'équation $f(x)=0$ admet une unique solution $\alpha$ dans $\R$.
  \item Montrer que $-1<\alpha<0$.
  \item À l'aide de la calculatrice, donner un encadrement de $\alpha$ d'amplitude
        $10^{-2}$.
  \item En déduire le signe de $f(x)$ suivant les valeurs de $x$.
\end{enumerate}
\end{exercice}

% ----------------------------------------------------------- Ex 11
\begin{exercice}[Dénombrer des solutions]{\niveauB}
Soit $f$ la fonction définie sur $\R$ par $f(x)=2x^3-3x^2-12x+1$.
\begin{enumerate}
  \item Calculer $f'(x)$ et étudier son signe.
  \item Dresser le tableau de variations complet de $f$ (limites comprises).
  \item Montrer que l'équation $f(x)=0$ admet exactement trois solutions
        $\alpha_1<\alpha_2<\alpha_3$.
  \item Déterminer, pour chacune d'elles, un encadrement par deux entiers consécutifs.
  \item Déterminer le nombre de solutions de l'équation $f(x)=10$, puis de
        l'équation $f(x)=-19$.
\end{enumerate}
\end{exercice}

% ----------------------------------------------------------- Ex 12
\begin{exercice}[Méthode de dichotomie]{\niveauB}
Soit $f$ la fonction définie sur $\intff{1}{2}$ par $f(x)=x^3+x-3$.
\begin{enumerate}
  \item Montrer que l'équation $f(x)=0$ admet une unique solution $\alpha$ dans
        $\intff{1}{2}$.
  \item Compléter le tableau suivant, qui applique la méthode de dichotomie
        (on indiquera le signe de $f(m)$) :
\begin{center}
\setstretch{1}
\renewcommand{\arraystretch}{1.6}
\newcolumntype{C}[1]{>{\centering\arraybackslash}m{#1}}
\begin{tabular}{|C{1.2cm}|C{1.8cm}|C{1.8cm}|C{2.2cm}|C{2.3cm}|C{3.4cm}|}
\hline
\textbf{Étape} & $\boldsymbol{a}$ & $\boldsymbol{b}$ & $\boldsymbol{m=\frac{a+b}{2}}$ &
  \textbf{Signe de} $\boldsymbol{f(m)}$ & \textbf{Nouvel intervalle} \\ \hline
1 & $1$ & $2$ & $1{,}5$ & & \\ \hline
2 & & & & & \\ \hline
3 & & & & & \\ \hline
4 & & & & & \\ \hline
\end{tabular}
\end{center}
  \item Compléter le programme Python suivant pour qu'il renvoie un encadrement de
        $\alpha$ d'amplitude inférieure ou égale à $p$.

\begin{minipage}{0.62\linewidth}
\begin{lstlisting}
def f(x):
    return x**3 + x - 3

def dichotomie(a, b, p):
    while ............ :
        m = ............
        if f(a) * f(m) <= 0:
            b = ............
        else:
            a = ............
    return a, b
\end{lstlisting}
\end{minipage}

  \item Combien d'étapes sont nécessaires pour obtenir un encadrement d'amplitude
        inférieure à $10^{-3}$ ?
\end{enumerate}
\begin{indication}
Après $n$ étapes, l'amplitude de l'intervalle est $\dfrac{2-1}{2^n}=\dfrac{1}{2^n}$.
\end{indication}
\end{exercice}

% =====================================================================
\section{Vers le bac}
% =====================================================================

% ----------------------------------------------------------- Ex 13
\begin{exercice}[TVI sur un intervalle ouvert]{\niveauC}
On souhaite étudier l'équation $(E)$ : $\e^{x}=\dfrac{1}{x}$ sur
$\intoo{0}{+\infty}$.

On considère la fonction $g$ définie sur $\intoo{0}{+\infty}$ par
$g(x)=\e^{x}-\dfrac{1}{x}$.
\begin{enumerate}
  \item Déterminer $\limd{0}{g(x)}$ et $\limite{x\to+\infty}{g(x)}$.
  \item Calculer $g'(x)$ et en déduire les variations de $g$.
  \item Montrer que l'équation $(E)$ admet une unique solution $\alpha$ dans
        $\intoo{0}{+\infty}$.
  \item Montrer que $0{,}5<\alpha<0{,}6$, puis donner un encadrement de $\alpha$
        d'amplitude $10^{-2}$.
  \item Montrer que $\e^{-\alpha}=\alpha$.
\end{enumerate}
\end{exercice}

% ----------------------------------------------------------- Ex 14
\begin{exercice}[Problème : fonction auxiliaire]{\niveauC}
\textbf{Partie A -- Étude d'une fonction auxiliaire}

Soit $g$ la fonction définie sur $\R$ par $g(x)=\e^{x}-x-2$.
\begin{enumerate}
  \item Déterminer la limite de $g$ en $-\infty$.
  \item En écrivant $g(x)=x\left(\dfrac{\e^{x}}{x}-1-\dfrac{2}{x}\right)$ pour
        $x\neq 0$, déterminer la limite de $g$ en $+\infty$.
  \item Étudier les variations de $g$ et dresser son tableau de variations.
  \item Montrer que l'équation $g(x)=0$ admet exactement deux solutions $\alpha$ et
        $\beta$, avec $\alpha<0<\beta$.
  \item Montrer que $-2<\alpha<-1$ et $1<\beta<2$, puis donner un encadrement de
        $\beta$ d'amplitude $10^{-1}$.
  \item En déduire le signe de $g(x)$ suivant les valeurs de $x$.
\end{enumerate}

\textbf{Partie B -- Étude de la fonction principale}

Soit $f$ la fonction définie sur $\R$ par $f(x)=\e^{x}-\dfrac{x^2}{2}-2x$.
\begin{enumerate}
  \item Montrer que, pour tout réel $x$, $f'(x)=g(x)$.
  \item En déduire le tableau de variations de $f$ (sans les limites).
  \item En utilisant l'égalité $\e^{\beta}=\beta+2$, montrer que
        $f(\beta)=-\dfrac{\beta^2}{2}-\beta+2$.
\end{enumerate}
\end{exercice}

% ----------------------------------------------------------- Ex 15
\begin{exercice}[Approfondissement : deux résultats classiques]{\niveauC}
\textbf{Partie A -- Théorème du point fixe}

Soit $f$ une fonction continue sur $\intff{0}{1}$ telle que, pour tout
$x\in\intff{0}{1}$, $0\leqslant f(x)\leqslant 1$.

On pose, pour tout $x\in\intff{0}{1}$, $\varphi(x)=f(x)-x$.
\begin{enumerate}
  \item Justifier que $\varphi$ est continue sur $\intff{0}{1}$ et déterminer le
        signe de $\varphi(0)$ et de $\varphi(1)$.
  \item En déduire qu'il existe un réel $c\in\intff{0}{1}$ tel que $f(c)=c$.
  \item Interpréter graphiquement ce résultat.
  \item \textbf{Application :} montrer que l'équation $\cos x=x$ admet au moins une
        solution dans $\intff{0}{1}$.
\end{enumerate}

\textbf{Partie B -- Théorème de la corde}

Soit $f$ une fonction continue sur $\intff{0}{1}$ telle que $f(0)=f(1)$.

On pose, pour tout $x\in\intff{0}{\frac{1}{2}}$,
$h(x)=f\!\left(x+\frac{1}{2}\right)-f(x)$.
\begin{enumerate}
  \item Justifier que $h$ est continue sur $\intff{0}{\frac{1}{2}}$.
  \item Calculer $h(0)+h\!\left(\frac{1}{2}\right)$.
  \item En déduire qu'il existe un réel $c\in\intff{0}{\frac{1}{2}}$ tel que
        $f\!\left(c+\frac{1}{2}\right)=f(c)$, puis interpréter graphiquement.
\end{enumerate}
\end{exercice}

\end{document}
